\section{由流方程生成的变换}

场空间中的流通过无穷小变换逐步构建出场变换。
与整体变换相比，无穷小变换通常更易于处理，
因为它们只涉及当前时刻的场。
此外，所生成的变换的可微性与可逆性
自动得到保证。


\subsection{场空间中的流}

一个无穷小场变换，
\begin{equation}
  U\to U+\eps Z(U)U+\rmO(\eps^2),
\end{equation}
由取值于 $\Lie$ 的链接场 $[Z(U)](x,\mu)$ 生成。
此类变换的连续复合
等价于对某个虚构时间 $t$ 积分流方程
\begin{equation}
  \dot{U}_t=Z_t(U_t)U_t.
\end{equation}
流生成元的一个简单选择是
\begin{align}
  [Z_t(U)]^a(x,\mu)&=\du^a_{x,\mu}\Wl_0(U),
  \\[1.5ex]
  \Wl_0(U)&=\sum_{x,\mu\neq\nu}\,\tr\{U(x,\mu,\nu)\},
\end{align}
其中 $U(x,\mu,\nu)$ 表示点 $x$ 处 $(\mu,\nu)$ 方向的方格（plaquette）圈。
下文中，这一流将被称为``威尔逊流''。

若 $Z_t(U)$ 是 $t$ 和 $U$ 的可微函数，则对任意给定的初值 $U_0=V$
以及所有 $t\in(-\infty,\infty)$，
流方程 (3.2) 都有唯一解 $U_t$。
而且，解关于 $t$ 和 $V$ 可微
（这些论断的证明可参见例如文献~\cite{ArnoldI}）。
应当强调的是，
解对所有时间都存在这一点并非平凡，
只有在不再附加其他假设、
而场流形为紧致时才能得到保证。


\subsection{积分后的变换}

在固定的 $t$ 处，场 $U_t$ 是初始场 $V$ 的一个良好定义的函数。
通过对流方程积分，便得到场空间的一个可微变换
$V\to U_t=\transt(V)$。
该变换是可逆的且其逆可微，
因为流方程可以从 $t$ 向 0 反向积分。
此外，
由于雅可比行列式 $\det\transts(V)$
在 $t=0$ 时等于 1 且在任何时刻都不经过零，
该变换还是保持定向的，
从而满足场空间可接受映射的全部要求。

雅可比行列式有一个有用的紧凑表达式，
它从下列方程出发得到：
\begin{align}
  {\rmd\over\rmd t}\ln\det\transts(V)
  &=\Tr\bigl\{\dot{\trans}_{t,*}(V)\transts(V)^{-1}\bigr\},
  \\[2.0ex]
  [\dot{\trans}_{t,*}(V)](x,\mu;y,\nu)^{ab}
  ={}&-2\,\tr\Bigl\{\hat{\du}^b_{y,\nu}\bigl\{[Z_t(U_t)](x,\mu)U_t(x,\mu)\bigr\}
  U_t(x,\mu)^{-1}T^a
  \notag\\[0.5ex]
  &\hspace{3.2em}-\hat{\du}^b_{y,\nu}U_t(x,\mu)U_t(x,\mu)^{-1}
  [Z_t(U_t)](x,\mu)T^a\Bigr\}.
\end{align}
注意到
\begin{equation}
  \hat{\du}^b_{y,\nu}=\sum_{x,\mu}[\transs(V)](x,\mu;y,\nu)^{ab}
  \du^a_{x,\mu},
\end{equation}
再经过几行代数运算便得到公式
\begin{equation}
  \ln\det\transts(V)=\int_0^t\rmd s\,\sum_{x,\mu}
  \bigl\{\du^a_{x,\mu}[Z_s(U)]^a(x,\mu)\bigr\}_{U=U_s}.
\end{equation}
例如在威尔逊流的情形，
雅可比行列式对场 $V$ 的作用量的贡献
\begin{equation}
  \ln\det\transts(V)=-\frac{16}{3}\int_0^t\rmd s\,\Wl_0(U_s)
\end{equation}
正比于 Wilson 方块作用量沿流的积分。
